\section{Systemsicherheit}
\subsection{Malware (\blue{MAL}ixious Soft\blue{WARE})}
\fatsf{Definition:} Schädlicher Code des Angreifers, der auf der Maschine des Opfers ausgeführt wird\\
\fatsf{Ziel:} Vielseitige Angriffe möglich, u.a.
\begin{itemize}
    \item Löschen von Dateien
    \item Senden von Spam Nachrichten von Opfermaschine
    \item Denial of Service Angriff als Teil eines Botnets
    \item Harddisk Verschlüsselung zur Erpressung von Lösegeld
    \item Private Informationen extrahieren (Passwörter etc.)
    \item Physikalische Zerstörung des Computers
\end{itemize}
\subsubsection*{Arten von Malware}
\begin{figure}[h]
    \centering
    \includegraphics[width=1\textwidth]{Thema_8_Images/Arten von Malware.jpeg}
\end{figure}
\subsubsection*{Spezialfälle}
\begin{itemize}
    \item \blue{Ransomware:} Trojaner, der Dateien verschlüsselt und nur gegen \enquote{Lösegeld}-Zahlung wieder freigibt (oder auch nicht)
    \item \blue{Scareware:} Trojaner, der z.B. als Virenscanner vermeintlich gefundene Viren gegen Bezahlung löschen soll
\end{itemize}
\begin{defBox}
    \fatsf{Zero-Day-Exploits:} Malware nutzt häufig neue Schwachstellen, bei der Entwickler keine Zeit hat (zero days), um Security Patch bereitzustellen
\end{defBox}
\subsubsection*{Malware als Spionage Tool}
\blue{Flame malware}-Angriff auf den Microsoft Windows Update Mechanismus nutzte \red{MD5} Kollisionen aus, um Zertifikat von Microsoft zu fälschen
\begin{figure}[h]
    \centering
    \includegraphics[width=0.6\textwidth]{Thema_8_Images/Malware als spionage Tool.jpeg}
\end{figure}
\subsubsection*{Malware verbreiten}
\begin{itemize}
    \item \blue{klassisch:} als E-Mail-Anhang oder während einer Programm-Installation
    \begin{itemize}
        \item Gegenmaßnahmen: \red{nur vertrauenswürdige Quellen} akzeptieren
    \end{itemize}
    \item \blue{drive-by-download:} beim Besuch einer Web-Seite (via Java, Javascript, Flash, ...)
    \begin{itemize}
        \item Gegenmaßnahmen: \red{Ausführung von Skripten deaktivieren, Sandboxing}
    \end{itemize}
    \item \blue{physisch:} durch infizierte Tokens
    \begin{itemize}
        \item Gegenmaßnahmen: \red{nur vertrauenswürdige Quellen} akzeptieren
    \end{itemize}
\end{itemize}
\begin{figure}[h]
    \centering
    \includegraphics[width=0.9\textwidth]{Thema_8_Images/Malware: Verbreiten.jpeg}
\end{figure}
\subsubsection*{Würmer}
\begin{itemize}
    \item Programme, die sich automatisch \red{über das Netz selbst verbreiten}
    \begin{itemize}
        \item \fatsf{Mit oder sogar ganz ohne Nutzerinteraktion}, daher auch die Klassifikation \enquote{wormable} für Schwachstellen, die komplett automatische Verbreitung erlauben
    \end{itemize}
    \item \blue{Beispiel:} Sasser (2004)
    \begin{itemize}
        \item Nutzte eine Schwachstelle (Buffer Overflow) im Windows 2000/XP Dienst LSASS aus, um sich selbst auf Zielsystem zu kopieren
        \begin{itemize}
            \item Patch zur Schließung der Schwachstelle war schon veröffentlicht, aber nicht überall installiert
        \end{itemize}
        \item Auf dem Zielsystem angekommen scannt der IP Bereiche nach neuen verwundbaren Zielen
        \item Schadfunktion: Herunterfahren des Computers in unregelmäßigen Abständen
    \end{itemize}
\end{itemize}
\subsubsection*{Viren}
\begin{itemize}
    \item Programme, die ihre \red{eigenen Code in \enquote{Wirtsprogramme}} einfügen und sich bei Nutzerinteraktion verbreiten
    \item Bekannte Virentypen:
    \begin{itemize}
        \item Executable infectors
        \item Polymorphic viruses
        \item Metamorphic viruses
        \item ...
    \end{itemize}
\end{itemize}
\subsubsection*{Viren: Executable Infectors}
\begin{itemize}
    \item Virus infiziert ein ausführbares Programm
    \item Virencode wird im Wirtsprogramm eingefügt und Kontrollfluss entsprechend adaptiert
    \item Aktion des Virus (\enquote{payload}) kann von externen Faktoren abhängig sein (z.B. Zeit, Netzwerkadresse, ...)
    \begin{infoBox}
        Ein \fatsf{Payload} ist die eigentliche \fatsf{Schadfunktion} eines Virus, Wurms oder einer anderen Malware.\\
        \fatsf{Vergleich:} Der Virus ist das \enquote{Transportmittel}, der Payload ist die \enquote{gefährliche Fracht}
    \end{infoBox}
    \begin{figure}[h]
        \centering
        \includegraphics[width=1\textwidth]{Thema_8_Images/Viren: Executable Infectors.jpeg}
    \end{figure}
    \item Erkennung einfach anhand des Virus Codes $\rightarrow$ statische Virussignatur
\end{itemize}
\subsubsection*{Statische Virussignatur}
\begin{itemize}
    \item Berechnung eines Algorithmus oder Hashes, um Virus eindeutig zu identifizieren
    \item Format: <Virus CRC16/CRC32/Hash Value> | <Virus Name>
    \begin{itemize}
        \item 0095C3A4|STONED.LESZOP.A
        \item 0086C7BE|STONED.MARCH6.A
    \end{itemize}
\end{itemize}
\subsubsection*{Viren: Polymorphe Viren}
\begin{itemize}
    \item Virus ist (bis auf kurze Entschlüsselungsroutine) \red{verschlüsselt}
    \begin{figure}[h]
        \centering
        \includegraphics[width=0.8\textwidth]{Thema_8_Images/Viren: Polymorphe Viren.jpeg}
    \end{figure}
    \item Code für die Entschlüsselung besagt: \enquote{Entschlüssele verschlüsselten Viruscode mit Schlüssel und führe den Virus Code aus}
\end{itemize}
\begin{itemize}
    \item Virus verändert Erscheinungsbild mit jeder Replikation
    \begin{figure}[h]
        \centering
        \includegraphics[width=0.8\textwidth]{Thema_8_Images/Viren: Polymorphe Viren 2.jpeg}
    \end{figure}
    \item Beide Viren nutzen \fatsf{unterschiedliche} Schlüssel $\rightarrow$ Alles bis auf den kurzen Code für die Entschlüsselung sieht \enquote{zufällig} aus
    \item \red{Erkennung deutlich schwieriger} (großer Teil ist zufälliger String) $\rightarrow$ Erkennung fokussiert auf Entschlüsselungsroutine
\end{itemize}
\subsubsection*{Viren: Metamorphe Viren}
Jedes mal wenn sich der Virus repliziert, erzeugt er eine semantisch unterschiedliche Version des Virus Codes\\
$\rightarrow$ Gleiche High-level Aktionen aber mit kleinen unbedeutenden Abweichungen\\
$\rightarrow$ Beispiele für äquivalenten Code
\begin{itemize}
    \item add 0 to register
    \item subtract 0 from register
    \item xor 0 with register
    \item no-op
    \item Veränderung von conditional statements (if/else)
    \item Ersetzen eines Algorithmus durch anderen (z.B. mergesort durch quicksort)
\end{itemize}
\begin{figure}[h]
    \centering
    \includegraphics[width=1\textwidth]{Thema_8_Images/Viren: Metamorphe Viren.jpeg}
\end{figure}
\begin{itemize}
    \item \red{Rwriter Code} besagt: \enquote{Erzeuge semantisch unterschiedlichen Code des Virus Codes}
    \item \red{Achtung:} Rewriter Code kann selbst modifiziert werden
\end{itemize}
\subsubsection*{Metamorphe Viren: Gegenmaßnahmen}
Erkennung anhand des Verhaltens:
\begin{itemize}
    \item Anstelle des Syntax muss das Verhalten des Virus analysiert werden
    \item Untersuchung der Aktionen, die der Virus ausführt
    \item Antiviren Firmen suchen nach Verhaltensmustern, um eine Virussignatur zu erstellen
\end{itemize}
Strategien eines Virus, um Erkennung zu verhindern:
\begin{itemize}
    \item Bösartiger Viruscode wird erst nach langer Wartezeit ausgeführt
    \item Erkenne Codeanalyse und stoppe die Ausführung des Virus Codes
\end{itemize}
$\rightarrow$ Katz- und Mausspiel zwischen Viren- und Antivirenentwicklern
\subsubsection*{Trojaner}
\begin{itemize}
    \item Programm, das neben bekannter Funktion \red{auch \enquote{unerwünschte} Nebeneffekte} hat
    \item Beispiel: Sunburst (2020)
    \begin{itemize}
        \item Kontroller über Systeme von Behörden weltweit (Dateien ableiten, Systemdienste deaktivieren)
        \item Aktiv seit März, gefunden im Dezember
        \item Einfallstor: Updateroutine der Software Orion (IT Administrationssoftware)
    \end{itemize}
    \item Thompson's Compiler (1984) - eher theoretisch interessant
    \begin{itemize}
        \item Modifiziere C Compiler so, dass bei Übersetzung von Tools wie Login-Modul bekannte Passwörter akzeptiert werden
        \item Sobald Compiler eine neue Version seines eigenen Codes übersetzt, wird der gleiche Schadcode eingefügt\\
        $\rightarrow$ Schwierig zu finden, da modifizierter Quellcode gelöscht und nur der modifizierter Compiler bereitgestellt wird
    \end{itemize}
\end{itemize}
\subsubsection*{Ransomware}
\begin{itemize}
    \item Bezeichnet eine \red{Schadfunktion}, keinen Verbreitungsweg
    \begin{itemize}
        \item Bsp. Wannacry: Ransomware und gleichzeitig Wurm
    \end{itemize}
    \item \red{Verschlüsselt nahezu alle Dateien auf den PC} (keine Systemdateien)
    \item z.B. hybride Verschlüsselung:
    \begin{itemize}
        \item jede Datei wird mit einem neuen symmetrischen Schlüssel verschlüsselt
        \item dieser wird mit dem public key des Angreifers verschlüsselt und an die Datei angehängt
    \end{itemize}
    \item Fordern Lösegeld (\enquote{ransom}) für Zugriff auf den privaten Schlüssel (häufig in Kryptowährungen)
    \item Meist ein Schlüsselpaar pro Opfer auf dem Server gespeichert
\end{itemize}
\subsection{Schutz vor viren}
\subsubsection*{Methoden der Antivirenprogramme}
\fatsf{Reaktiv}
\begin{itemize}
    \item Systemscan: sucht auf Festplatte nach verdächtigen Dateien (z.B. anhand statischer \fatsf{Virensignaturdatenbank})
    \item Real-Time Detection (On-Access Scanning): überprüft beim Speicher, Öffnen, Ausführen, ...
    \begin{itemize}
        \item z.B., schützt Ihr System vor Malware-Bedrohungen, während Sie im Internet surfen, Ihren Posteingang überprüfen oder Dateien herunterladen
    \end{itemize}
\end{itemize}
\fatsf{Proaktiv}
\begin{itemize}
    \item Activily Monitoring: verwendet Strategien zur Überwachung des \fatsf{virentypisches Verhalten}
    \begin{itemize}
        \item z.B. Replikation in AutoStart, auffällige DNS-Anfragen, ...
    \end{itemize}
    \item Heuristische Erkennungsalgorithmen zur Identifizierung verdächtiger Dateien und Codeteile, die noch nicht als Malware klassifiziert werden
\end{itemize}
\subsubsection*{Reaktives Scannen}
\begin{figure}[h]
    \centering
    \includegraphics[width=1\textwidth]{Thema_8_Images/reaktives Scannen.jpeg}
\end{figure}
\begin{itemize}
    \item Prüfe, ob Signatur bekannt
    \item Prüfe auch auf Mutationen, 
    \begin{itemize}
        \item z.B. wenn Virus verschlüsselt, und sich erst in Sandbox entschlüsseln lässt
    \end{itemize}
\end{itemize}
\subsubsection*{EICAR-Test für Virenschutz}
\begin{itemize}
    \item Einfaches Testen von Antiviren-Programmen durch (harmlose) Testdatei von \red{EICAR} (European Institute for Computer Anti-Virus Research e.V.):
\end{itemize}
\begin{figure}[h]
    \centering
    \includegraphics[width=1\textwidth]{Thema_8_Images/EICAR-Test.jpeg}
\end{figure}
\subsubsection*{Effektivität von Virenscannern?}
\enquote{People thought that virus protected them, but we can \fatsf{never} block all viruses}
\begin{itemize}
    \item Erfolgsstatistik von AV-Software variiert zwischen 40\% und 99.9\%
    \item Tests beziehen sich aber meistens auf bekannte Massen-Malware
    \item Erfolgsquote bei neuen gezielten Virenangriffen unklar
    \item Was ist mit Verzögerung bei Signatur-Updates?\\
    $\rightarrow$ Verhaltensbasierte Malware-Erkennung
\end{itemize}
\subsection{Ausführung von C Programmen}
\subsubsection*{Call: Compiler, Assembler, Linker, Loader}
\fatsf{4 Schritte zur Ausführung eines C-Programms:}
\begin{itemize}
    \item \fatsf{Compiler:} Konvertiert C-Code in Assembler-Code (RISC-V, x86)
    \item \fatsf{Assembler:} Konvertiert Assembler-Code in Maschinencode (raw bits)
    \item \fatsf{Linker:} Löst Abhängigkeiten von externen Bibliotheken auf (brauchen wir hier nicht)
    \item \fatsf{Loader:} Einrichtung von Speicherplatz und Ausführung des Maschinencodes
\end{itemize}
\begin{figure}[h]
    \centering
    \includegraphics[width=1\textwidth]{Thema_8_Images/Call.jpeg}
\end{figure}
\subsubsection*{Memory Layout}
\fatsf{Loader} weist das Betriebssystem an, dem Programm einen großen Speicherbereich zur Ausfürhung zuzuweisen
\begin{itemize}
    \item Jedes \fatsf{Byte} des Speichers hat eine eindeutige \fatsf{Adresse}
    \item Bei einem \fatsf{32-Bit-System} hat der Speicher \fatsf{32-Bit-Adressen}
    \item Jede \fatsf{Adresse} bezieht sich auf ein \fatsf{Byte}, d.h. wir haben $2^{32}$ Bytes Speicherplatz (32-Bit-System)
\end{itemize}
Häufige Darstellung des Speichers als \fatsf{Tabelle:}
\begin{itemize}
    \item In jeder \fatsf{Zelle} steht ein \fatsf{Byte}
    \item Die Zellen einer Zeile können zusammengefasst werden zu einem \fatsf{Wort (4 Bytes = 1 Wort)}
\end{itemize}
\begin{figure}[h]
    \centering
    \includegraphics[width=0.3\textwidth]{Thema_8_Images/Memory Layout.jpeg}
\end{figure}
\subsection{X86 Architektur}
\subsubsection*{X86 Memory Layout}
\fatsf{Speichernutzung bei Ausführung eines Programms}
\begin{itemize}
    \item \fatsf{Code:}
    \begin{itemize}
        \item Der Programmcode selbst (acuh \enquote{Text} genannt) in Maschinencode
    \end{itemize}
    \item \fatsf{Daten:}
    \begin{itemize}
        \item Statische Variablen/Konstanten, die beim Start des Programms zugewiesen werden
    \end{itemize}
    \item \fatsf{Heap:}
    \begin{itemize}
        \item Dynamisch zugewiesener Speicher mit \fatsf{malloc} und \fatsf{free} 
    \end{itemize}
    \item \fatsf{Stack:}
    \begin{itemize}
        \item Lokale Variablen und Stack Frames
        \item Bei tieferen Funktionsaufrufen, wächst er nach unten
    \end{itemize}
\end{itemize}
\subsubsection*{Stack Frame}
\begin{itemize}
    \item Wenn der Code eine Funktion aufruft, wird auf dem Stack Platz für lokale Variablen geschaffen
    \begin{itemize}
        \item Dieser Platz wird als \fatsf{Stack Frame} für die Funktion bezeichnet
        \item Der Stack Frame verschwindet, sobald die Funktion zurückkehrt
    \end{itemize}
    \item Der Stack startet an höherer Adresse
    \item Jedes Mal, wenn Code eine Funktion aufruft, \fatsf{wächst der Stack nach unten} und schafft zusätzlichen Platz
    \begin{itemize}
        \item Hinweis: Daten auf dem Stack, z.B. eine Zeichenkette (String), werden nach wie vor von der niedrigsten Adresse bis zur höchsten Adresse gespeichert
        \item Das \enquote{Herunterwachsen} geschieht nur, wenn zusätzlicher Speicher zugewiesen werden muss
    \end{itemize}
\end{itemize}
\begin{figure}[h]
    \centering
    \includegraphics[width=0.3\textwidth]{Thema_8_Images/Stack Frame.jpeg}
\end{figure}
\subsubsection*{Register}
Spezielle \fatsf{Speicherbereiche} auf der \fatsf{CPU}, die jeweils \fatsf{4 Bytes} speichern können
\begin{itemize}
    \item \fatsf{Adressierung} unterscheidet sich zu Memory Layout
    \begin{itemize}
        \item \fatsf{Speicher:} Adressen als 32-Bit-Zahlen
        \item \red{Register werden mit Namen bezeichnet} (EBP, ESP, EIP, ...), nicht mit Adressen
    \end{itemize}
    \item Um Überblick über aktuellen Stack-Frame zu behalten, speichern wir zwei Zeiger in Register:
    \begin{itemize}
        \item \fatsf{EBP}-Register (\enquote{Basiszeiger}) zeigt auf den Anfang des aktuellen Stack Frames
        \item \fatsf{ESP}-Register (\enquote{Stackzeiger}) zeigt auf das untere Ende des aktuellen Stack-Frame
        \item \fatsf{EIP}-Register (\enquote{Befehlszeiger}) speichert die Adresse des gerade ausgeführten Maschinenbefehls im Code Segment
    \end{itemize}
\end{itemize}
\begin{figure}[h]
    \centering
    \includegraphics[width=0.3\textwidth]{Thema_8_Images/Register.jpeg}
\end{figure}
\subsubsection*{Pushing und Popping}
\begin{itemize}
    \item Die \fatsf{Push}-Anweisung fügt dem Stack ein Element aus dem Register hinzu
    \begin{itemize}
        \item Dekrementiere den ESP, um mehr Speicher auf dem Stack zuzuweisen
        \item Speichern des neuen Wertes auf dem niedrigsten Wert des Stacks
    \end{itemize}
\end{itemize}
\clearpage
\begin{itemize}
    \begin{figure}[h]
        \centering
        \includegraphics[width=0.75\textwidth]{Thema_8_Images/pushing .jpeg}
    \end{figure}
    \item Die \fatsf{Pop}-Anweisung entfernt ein Element vom Stack
    \begin{itemize}
        \item Laden des Wertes vom niedrigsten Wert auf dem Stack und speichern in einem Register
        \item ESP inkrementieren, um den Speicher auf dem Stack freizugeben
    \end{itemize}
    \begin{figure}[h]
    \centering
    \includegraphics[width=0.75\textwidth]{Thema_8_Images/popping.jpeg}
\end{figure}
\end{itemize}
\subsection{X86 Funktionsaufruf}
\subsubsection*{Aufruf einer Funktion}
\begin{figure}[h]
    \centering
    \includegraphics[width=0.75\textwidth]{Thema_8_Images/Aufruf einer Funktion.jpeg}
\end{figure}
\subsubsection*{Funktionsaufruf in X86}
\begin{itemize}
    \item Beim Aufruf einer Funktion müssen ESP und EBP verschoben werden, um einen neuen Stack Frame zu erstellen
    \item EIP wird in den Code der aufgerufenen Funktion verschoben
    \item Bei der Rückkehr aus einer Funktion müssen ESP, EBP und EIP auf ihre alten Werte zurückkehren
\end{itemize}
\begin{figure}[h]
    \centering
    \includegraphics[width=1\textwidth]{Thema_8_Images/Funktionsaufruf in X86.jpeg}
\end{figure}
\subsubsection*{Wiederholung: Stack, Register}
\begin{itemize}
    \item Während Funktionsaufruf wird auf dem Stack Platz für lokale Variablen geschaffen (Stack Frame)
    \item Dieser Platz wird gelöscht, sobald die Funktion zurückkehrt
    \item Der Stack beginnt bei höheren Adressen und wächst nach unten
    \item Der Codeabschnitt enthält in Raw Bit, die Assembleranweisungen der Funktionen
    \item Wir speichern zwei Zeiger, die den Umfang des aktuellen Stack Frames definieren
    \item EBP wird für den oberen Teil des Stack Frames verwendet
    \item ESP für den unteren Teil des Stack Frames 
    \begin{itemize}
        \item Alles unter ESP ist undefiniert und wird ignoriert
        \item Beim Schreiben in Stack muss ESP angepasst werden, so dass er auf niedrigsten Wert referenziert
    \end{itemize}
    \item EIP zeigt auf den aktuellen Befehl in der Programmausführung
\end{itemize}
\subsubsection*{0. Initialer Stack}
\begin{itemize}
    \item Im Stack existiert ein Stack Frame für Funktion \fatsf{main}
    \item Jeder Stack Frame muss Platz für lokale Variablen haben und wächst nach unten (hier keine lokalen Variablen)
\end{itemize}
\fatsf{Merke:} Alle Werte der Register EBP, ESP und EIP müssen auf dem Stack abgelegt werden bei einem Funktionsaufruf, da diese sonst verloren gehen
\subsubsection*{1. Argumente der aufgerufenen Funktion}
\begin{itemize}
    \item Die Argumente werden dem Stack in umgekehrter Reihenfolge hinzugefügt
    \item ESP muss auf den untersten Wert des Stack Frames für die Funktion \fatsf{main} verweisen
\end{itemize}
\subsubsection*{2. Zwischenspeicherung EIP}
\begin{itemize}
    \item Als nächstes wird der aktuelle Wert von EIP auf den Stack gelegt
    \begin{itemize}
        \item Dies ermöglicht das Programm fortzusetzen nachdem die Funktion \fatsf{foo} zurückgekehrt ist
    \end{itemize}
    \item Wird manchmal auch als RIP (Return Instruction Pointer) bezeichnet
    \begin{itemize}
        \item Er zeigt, wohin wir im Programm in der Funktion \fatsf{main} als nächstes gehen müssen
    \end{itemize}
\end{itemize}
\subsubsection*{3. Zwischenspeicherung EBP}
\begin{itemize}
    \item Als nächstes wird der aktuelle Wert von EBP auf den Stack gelegt
    \begin{itemize}
        \item Dadurch können wir den obersten Punkt des vorherigen Stack Frames wiederherstellen, wenn wir zurückkehren
    \end{itemize}
\end{itemize}
\subsubsection*{4. Einstellen des Stack Frames}
\begin{itemize}
    \item Um den Stack Frame anzupassen, müssen wir alle drei Register aktualisieren
    \item Gefahrlos möglich da die alten Werte von EBP und EIP gespeichert wurden
    \item Altes ESP ist immer das untere Ende des Stacks, daher besteht keine Notwendigkeit, es zu speichern
\end{itemize}
\subsubsection*{5. Führe die Funktion aus}
\begin{itemize}
    \item Stack Frame bereit, die Anweisungen der Funktion zu verarbeiten
    \item Alle lokalen Variablen können auf Stack verschoben werden
    \item Durch aktuellen Wert von EBP können wir auch auf Variablen außerhalb des aktuellen Stack Frames zugreifen (z.B. Argumente)
\end{itemize}
\subsubsection*{6. Wiederherstellen}
\begin{itemize}
    \item Nach Beendigung der Funktion setzen wir alle drei Register wieder an ihren alten Wert zurück
    \item Verwende die in RIP und SFP gespeicherten Adressen, um EIP und EBP auf ihre alten Werte zurückzusetzen (durch Aufruf von pop)
    \item ESP wird auf das Ende des Stack Frames von \fatsf{main} zurückgesetzt 
    \begin{itemize}
        \item Beachte: Die Werte, die wir auf Stack geschoben haben sind immer noch da (wir überschreiben sie nicht, um Zeit zu sparen)
        \item Sie befinden sich unter ESP, so dass \fatsf{main} nicht mehr darauf zugreifen kann
    \end{itemize}
\end{itemize}
\clearpage
\begin{figure}[h]
    \centering
    \includegraphics[width=0.5\textwidth]{Thema_8_Images/Stack.jpeg}
\end{figure}
\subsection{Buffer Overflow}
\subsubsection*{2023 CWE TOP 25 most dangerous software weaknesses}
\begin{figure}[h]
    \centering
    \includegraphics[width=0.9\textwidth]{Thema_8_Images/Buffer Overflow.jpeg}
\end{figure}
\clearpage
\subsubsection*{Buffer Overflow}
\begin{itemize}
    \item Weit verbreitet insbesondere in C/C++
    \item C hat kein Konzept der Array Länge und sieht nur eine List von Bytes
    \begin{figure}[h]
        \centering
        \includegraphics[width=0.8\textwidth]{Thema_8_Images/Buffer Overflow 2.jpeg}
    \end{figure}
    \item Erstes bekanntes Beispiel: Internet Wurm Morris (1988)
    \item Vermeidung von Buffer Overflows: Kein Einsatz von C/C++\\
    $\rightarrow$ Ist jedoch nicht immer möglich
\end{itemize}
\subsubsection*{Anfälliger Code}
\begin{figure}[h]
    \centering
    \includegraphics[width=0.9\textwidth]{Thema_8_Images/Anfälliger Code .jpeg}
\end{figure}
\begin{infoBox}
    \begin{enumerate}
        \item \fatsf{Speicheranordnung im statischen Speicher:}\\
        Da name UNTERHALB von instrux im Speicher liegt, kann eine zu lange Eingabe instrux überschreiben
        \item \fatsf{Benutzer gibt eine lange Eingabe ein:}\\
        Wenn du mehr als 20 Zeichen eingibst, überschreibst du instrux:
        \begin{codeBlock}{}
            Gib deinen Namen ein: aaaaaaaaaaaaaaaaaaaaaaaaaa (mehr als 20 Zeichen)
        \end{codeBlock}
        \begin{itemize}
            \item Dadurch wird instrux überschrieben und könnte sich z.B. in \fatsf{bösartigen Code} verwandeln!
        \end{itemize}
    \end{enumerate}
\end{infoBox}
\begin{figure}[h]
    \centering
    \includegraphics[width=1\textwidth]{Thema_8_Images/Anfälliger Code 2.jpeg}
\end{figure}
\begin{figure}[h]
    \centering
    \includegraphics[width=1\textwidth]{Thema_8_Images/Anfälliger Code 3.jpeg}
\end{figure}
\begin{figure}[h]
    \centering
    \includegraphics[width=1\textwidth]{Thema_8_Images/Anfälliger Code 4.jpeg}
\end{figure}
\begin{figure}[h]
    \centering
    \includegraphics[width=1\textwidth]{Thema_8_Images/Anfälliger Code 5.jpeg}
\end{figure}
\clearpage
\subsection{Stack Smashing}
\begin{itemize}
    \item Die häufigste Art des Buffer Overflows
    \item Tritt im Stack Speicher auf
    \item Wiederholung:
    \begin{itemize}
        \item Was sind Werte auf dem Stack, die ein Angreifer überlaufen lassen kann?
        \begin{itemize}
            \item Lokale Variablen
            \item Funktionsargumente
            \item Gespeicherter Framezeiger (SFP):
            \item Rückgabe-Instruktion Zeiger (RIP): Springt auf diese Anweisung nach Rückkehr aus Funktion
        \end{itemize}
        \item Beispiel: Bei Rückkehr aus einer Funktion setzen wir EIP auf den im Stack abgelegten RIP Wert
        \begin{itemize}
            \item Wenn Angreifer RIP überschreibt, kann er an eine beliebige Adresse springen und dort Shellcode ausführen
        \end{itemize}
    \end{itemize}
\end{itemize}
\subsubsection*{Überschreiben des RIP}
\begin{figure}[h]
    \centering
    \includegraphics[width=1\textwidth]{Thema_8_Images/Überschreiben des RIP 1.jpeg}
\end{figure}
\clearpage
\begin{figure}[h]
    \centering
    \includegraphics[width=1\textwidth]{Thema_8_Images/Überschreiben des RIP 2.jpeg}
\end{figure}
\subsubsection*{Konstruktion von Exploits}
\begin{figure}[h]
    \centering
    \includegraphics[width=1\textwidth]{Thema_8_Images/Konstruktion von Exploits 1.jpeg}
\end{figure}
\clearpage
\begin{figure}[h]
    \centering
    \includegraphics[width=1\textwidth]{Thema_8_Images/Konstruktion von Exploits 2.jpeg}
\end{figure}
\subsubsection*{Beispiel: Remote Stack Smashing}
\begin{itemize}
    \item Overflow in Windows animated cursors (ANI) - \fatsf{LoadAniIcon()}
    \begin{itemize}
        \item Beschreibung:
        \begin{itemize}
            \item Animierte Cursordateien (.ani) enthalten animierte Grafiken für Symbole und Cursor
            \item Ein Stack Overflow tritt auf, da Microsoft Windows nicht die Größe des Headers einer fehlerhaften animierten Cursordateien überprüft
        \end{itemize}
        \item Auswirkungen:
        \begin{itemize}
            \item Ein Angreifer kann beliebigen Code ausführen oder einen Denial-of-Service-Angriff durchführen
        \end{itemize}
    \end{itemize}
    \item Buffer overflow in Symantec virus detection (Mai 2016)
    \begin{itemize}
        \item Buffer overflow beim Parsen von ausführbaren gepackten Dateien \\
        $\rightarrow$ Kernel Memory Korruption
        \item Häufigste Auswirkung: Erfolgreicher Angriff führt zu Systemabsturz oder Bluescreen of Death
    \end{itemize}
\end{itemize}
\clearpage
\subsection{Integer Overflow}
\subsubsection*{Signed/Unsigned Schwachstelle}
Ist folgender Code sicher?
\begin{figure}[h]
    \centering
    \includegraphics[width=0.85\textwidth]{Thema_8_Images/Signed-Unsigned Schwachstelle.jpeg}
\end{figure}
\begin{infoBox}[title=Was passiert bei Signed/Unsigned-Mischung?]
    Ein gefährliches Problem entsteht, wenn len ein \fatsf{signed int} ist, aber mit \fatsf{einem unsigned Wert} verglichen wird
    \begin{codeBlock}{}
        void func(int len) {
            if (len > (size_t)64)  // len wird zu unsigned gecastet!
                printf("len ist größer als 64\n");
        }
    \end{codeBlock}
    \begin{itemize}
        \item \fatsf{Gefahr:} Wenn len == -1, dann wird es zu 4294967295 auf 32-Bit-Systemen gecastet!
        \item \fatsf{Ergebnis:} -1 > 64 ist plötzlich WAHR!
    \end{itemize}
\end{infoBox}
\begin{figure}[h]
    \centering
    \includegraphics[width=0.9\textwidth]{Thema_8_Images/signed-unsigned schwachstelle 2.jpeg}
\end{figure}
\subsubsection*{Integer Overflows verhindern}
Von allen Integers müssen die Bereiche überprüft werden
\begin{codeBlock}{}
    void safe (size_t len, char *data){
        if (len > SIZE_MAX - 2) //Überprüfe ob Überlauf möglich
            return;
        char *buf = malloc(len + 2);
        if (!buf)
            return;
        memcpy(buf, data, len);
        buf[len] = '\n';
        buf[len + 1] = '\0';
    }
\end{codeBlock}
Sehr aufwendig und fehleranfällig!
\subsection{Gegenmaßnahmen}
\subsubsection*{Memory Safe Programming Languages}
\begin{itemize}
    \item Hauptgrund für Buffer Overflows: C/C++ ist nicht Memory Safe
    \item Lösung: Benutze Programmiersprache mit \enquote{Memory Safety} (z.B. Java, Python, Go, Rust, ...)
    \begin{itemize}
        \item Überprüfe reservierte Speichergröße und verhindere undefinierten Zugriff auf Speicherbereiche
    \end{itemize}
    \item Warum werden nicht nur Programmiesprachen mit \enquote{Memory Safety} verwendet?
    \begin{itemize}
        \item Effizienz: in vielen Anwendungen nicht relevant
        \item Legacy Code
    \end{itemize}
\end{itemize}
\subsubsection*{Programmieren}
\begin{itemize}
    \item Defensives Programmieren: Füge Gültigkeitsprüfungen ein wo immer möglich 
    \begin{itemize}
        \item Beispiel: Überprüfe, ob ein Zeiger null ist bevor er dereferenziert wird
    \end{itemize}
    \item \enquote{Gefährliche} Operationen wie \fatsf{gets/strcpy/strcat/sprintf} vermeiden
    \begin{itemize}
        \item Varianten fgets/strncpy/strncat/snprintf benutzen
        \item Aber das ist gar nicht so einfach, z.B. \fatsf{gets(buf) $\rightarrow$ fgets(buf, size, src)}
        \item Wobei \fatsf{size} so gesetzt werden sollte, dass nicht über \fatsf{buf} hinaus geschrieben werden kann
    \end{itemize}
    \item Verwendung von Preconditions, Postconditions und Invarianten, die der Code erfüllen muss
    \begin{itemize}
        \item Sehr aufwendig $\rightarrow$ daher kaum verwendet in der Praxis
    \end{itemize}
    \item Testen
    \begin{itemize}
        \item z.B. Fuzzing: Teste Programmcode auf zufälligen Eingaben
        \item Problem: Code Coverage
    \end{itemize}
\end{itemize}
\subsubsection*{Ausnutzen eines Memory Safety Schwachstelle}
\begin{enumerate}
    \item Finde eine Memory Safety Schwachstelle (z.B. Buffer Overflow)
    \item Schreibe Shellcode des Angreifers an eine bekannte Adresse des Speichers
    \begin{itemize}
        \item Randomisierung der Adressen im Speicherbereich
    \end{itemize}
    \item Überschreibe den RIP mit der Adresse des Shellcodes
    \begin{itemize}
        \item Gegenmaßnahme: Stack canaries
    \end{itemize}
    \item Springe aus der angegriffenen Funktion zurück
    \item Führe Shellcode des Angreifers aus
    \begin{itemize}
        \item Markiere Speicherbereiche als nicht ausführbar
    \end{itemize}
\end{enumerate}
\subsubsection*{Gegenmaßnahmen: Stack Canary}
\fatsf{Idee:} Füge einen zufälligen Wert zum Stack hinzu und überprüfe, ob er verändert wurde
\begin{itemize}
    \item Beim Start des Programms erzeuge einen zufälligen geheimen Wert (\enquote{\blue{canary}}) und speichere ihn in einer globalen Variable
    \item Beim Beginn der Ausführung der Funktion speichere \blue{canary} direkt unterhalb SFP/RIP
    \item Beim Beenden der Funktion überprüfe, ob \blue{canary} gleich dem Wert der globalen Variable ist\\
    $\rightarrow$ Angreifer kann RIP nur überschreiben, wenn er \blue{canary} zerstört\\
    $\rightarrow$ Erkenne Buffer Overflow Angriff $\rightarrow$ Führe das Programm zum Absturz
\end{itemize}
\begin{figure}[h]
    \centering
    \includegraphics[width=0.3\textwidth]{Thema_8_Images/Stack Canary.jpeg}
\end{figure}
Stack Canaries werden vom Compiler eingefügt
\begin{itemize}
    \item z.B. in Visual C++ im Modus /GS (GuardStack)
    \item Fügt automatisiert ein paar Befehle zu anfälligen Code hinzu $\rightarrow$ meistens nur geringer Effizienzverlust
\end{itemize}
Stack Canaries umgehen:
\begin{itemize}
    \item Extrahiere den Wert des Stack Canaries: z.B. durch andere Schwachstelle im Programm
    \item Verhindere das Überschreiben des Stack Canaries: z.B. bestimmte Angriffe erlauben Überschreiben beliebiger Speicherpositionen
    \item Raten des Stack Canaries: z.B. relativ einfach auf 32-bit Systemen
\end{itemize}
\subsubsection*{Gegenmaßnahmen: Nicht ausführbare Speicherbereiche}
Speicherbereich muss häufig nicht gleichzeitig ausführbar und beschreibbar sein
\begin{itemize}
    \item Code: Ausführbar aber nicht beschreibbar
    \item Stack, Heap: Beschreibbar aber nicht ausführbar
    \item \fatsf{DEP:} Data Execution Prevention (alias $W_{rite} \oplus$ eXecute) in Windows
    \begin{itemize}
        \item Teile des Speichers werden als nicht ausführbar markiert\\
        $\rightarrow$ Shellcode kann nicht mehr ausgeführt werden
    \end{itemize}
\end{itemize}
\fatsf{Probleme:}
\begin{itemize}
    \item Im Code Speicherbereich befinden sich bereits zahlreiche ausführbare Befehle
    \begin{itemize}
        \item Beispiel: \fatsf{Return-Oriented Programming (ROP)}
    \end{itemize}
\end{itemize}
\subsubsection*{Return-Oriented Programming}
\begin{infoBox}
    \fatsf{Return-Oriented Programming (ROP)} ist eine Technik, bei der \fatsf{bereits bestehende, ausführbare Codefragmente} (sogenannte \enquote{Gadgets}) genutzt werden, um den Codefluss \fatsf{umzuprogrammieren}. Diese Technik wird oft in \fatsf{Systemen verwendet}, bei denen DEP aktiviert ist
\end{infoBox}
\begin{figure}[h]
    \centering
    \includegraphics[width=0.67\textwidth]{Thema_8_Images/Return-Oriented Programming.jpeg}
\end{figure}
\begin{infoBox}
    \begin{enumerate}
        \item \fatsf{Angreifer sucht nach kurzen, wiederverwendbaren Code-Schnipseln} (Gadgets), die auf dem Stack oder im Speicher des Programms vorhanden sind
        \item Diese \fatsf{Gadgets} enthalten Rückkehrbefehle (z.B. RET-Befehle), die vom Angreifer genutzt werden, um den Programmfluss umzuleiten un die Kontrolle zu übernehmen
        \item Anstatt neuen Schadcode auszuführen, kombiniert der Angreifer einfach vorhandene Codefragmente in einer \fatsf{bestimmten Reihenfolge}, um die gewünschte Funktion zu erreichen
    \end{enumerate}
\end{infoBox}
\subsubsection*{Gegenmaßnahmen: Randomiesierung}
\fatsf{ASLR:} Adress Space Layout Randomization
\begin{itemize}
    \item Angriffe benötigen Kenntnis an welcher Stelle im speicher Schadcode platziert wurde
    \item Idee: Randomisiere Adressbereich eines Programms\\
    $\rightarrow$ Angreifer kann RIP nicht mehr so leicht korrekt setzen
\end{itemize}
\fatsf{ASLR} umgehen:
\begin{itemize}
    \item Spraying: Verbreite den Schadecode großflächig im Speicher
    \item Raten der Adresse: nur möglich auf 32-Bit Systemen
\end{itemize}
\subsection{Seitenkanalangriffe}
\subsubsection*{Angriffsvektoren}
\begin{figure}[h]
    \centering
    \includegraphics[width=0.9\textwidth]{Thema_8_Images/Angriffsvektoren.jpeg}
\end{figure}
\subsubsection*{Seitenkanalangriffe}
\fatsf{Idee:}
\begin{enumerate}
    \item \blue{Beobachte das Gerät bei Ausführung} des Algorithmus
    \begin{itemize}
        \item Messe die Zeit, die für Operationen benötigt werden
        \begin{itemize}
            \item Timing attacks
        \end{itemize}
        \item Messe den Stromverbrauch
        \begin{itemize}
            \item Power analysis attacks
        \end{itemize}
    \end{itemize}
    \item \blue{Füge Fehler in Berechnung ein} (z.B. mit einem Laserstrahl)
    \begin{itemize}
        \item Beobachte die Veränderungen im Verhalten
        \item Leite dadurch Informationen über Daten ab (z.B. über kryptographischen Schlüssel)
    \end{itemize}
\end{enumerate}
\subsubsection*{Beispiel: Passwort Check}
\begin{figure}[h]
    \centering
    \includegraphics[width=1\textwidth]{Thema_8_Images/Beispiel: Passwort Check.jpeg}
\end{figure}
Angriff:
\begin{itemize}
    \item Messe Laufzeit für verschiedene Eingaben
    \item Probiere erst alle möglichen Zeichen für die erste Position, wenn Laufzeit steigt, dann Zeichen korrekt
    \item Wiederhole für die folgende Position
\end{itemize}
\begin{figure}[h]
    \centering
    \includegraphics[width=1\textwidth]{Thema_8_Images/Beispiel: Passwort Check 2.jpeg}
\end{figure}
\begin{itemize}
    \item Erwartete Laufzeit für Angriff auf Passwort der Länge l mit m möglichen Zeichen:
    \begin{itemize}
        \item Bruteforce Angriff: $\frac{1}{2} \cdot m^l$
        \begin{itemize}
            \item Für Passworter mit 26 Zeichen der Länge 10 $\rightarrow$ $\approx$ 70.000.000.000.000 Tests
        \end{itemize}
        \item Timing Angriff: $\frac{1}{2} \cdot m \cdot l$
        \begin{itemize}
            \item Für Passwörter mit 26 Zeichen der Länge 10 $\rightarrow$ $\approx$ 130 Tests
        \end{itemize}
    \end{itemize}
    \item Laufzeit auch z.B. über Netzwerk messbar
    \item In der Praxis oft mehrere Messungen pro Eingabe notwendig, um Noise zu vermeiden
    \item Fundamentales Problem:
    \begin{itemize}
        \item \fatsf{Messbare Größe (Laufzeit, Stromverbrauch, ...) ist abhängig von geheimen Daten (Passwort, AES Schlüssel, ...)}
    \end{itemize}
\end{itemize}
\subsection{Mikroarchitektur-Angriff: Meltdown und Spectre}
\subsubsection*{Meltdown und Spectre}
\begin{itemize}
    \item Mikroarchitektur-Angriffe
    \item Entdeckt in 2017 von 4 verschiedenen Forschergruppen
    \item Veröffentlicht im Januar 2018
    \item Seitdem einige neue Varianten
\end{itemize}
\begin{figure}[h]
    \centering
    
\includegraphics[width=0.35\textwidth]{Thema_8_Images/Meltdown und Spectre.jpeg}
\end{figure}
\subsection{Caches und Cache Timing Angriffe}
\subsubsection*{Cache und Cache Timing}
\begin{figure}[h]
    \centering
    \includegraphics[width=0.8\textwidth]{Thema_8_Images/Cache und Cache Timing.jpeg}
\end{figure}
\begin{infoBox}
    \begin{enumerate}
        \item Der Angreifer führt einen Zugriff auf den Speicher durch
        \item Er misst, wie lange ein Lesezugriff dauert:
        \begin{itemize}
            \item \fatsf{Kurze Zeit:} Daten waren bereits im Cache (daher vorher verwendet)
            \item \fatsf{Lange Zeit:} Daten mussten aus dem Hauptspeicher geladen werden
        \end{itemize}
        \item Durch \fatsf{wiederholtes Testen} kann der Angreifer herausfinden, welche Speicherbereiche benutzt werden
    \end{enumerate}
\end{infoBox}
\subsubsection*{Cache Timing in der Praxis}
\begin{figure}[h]
    \centering
    \includegraphics[width=1\textwidth]{Thema_8_Images/Cache Timing in der Praxis.jpeg}
\end{figure}
\subsubsection*{Ziel des Angreifers}
\begin{figure}[h]
    \centering
    \includegraphics[width=1\textwidth]{Thema_8_Images/Ziel des Angreifers.jpeg}
\end{figure}
\begin{infoBox}
    \begin{itemize}
        \item \fatsf{Privilegiert (Kernel-Mode, Ring 0):}
        \begin{itemize}
            \item Hier läuft das \fatsf{Betriebssystem} (Kernel, Treiber, bestimmte Hardwarezugriffe)
            \item Hat Zugriff auf \fatsf{alle Instruktionen} und \fatsf{gesamten Speicher}
            \item Darf I/O-Hardware direkt ansprechen
            \item Beispiel: Kernelcode liest oder schreibt eine Prozess-Tabelle
        \end{itemize}
        \item \fatsf{Nicht-privilegiert (User-Mode, Ring 3):}
        \begin{itemize}
            \item Hier laufen \fatsf{normale Programme} (Browser, Spiele, Texteditoren)
            \item Zugriff \fatsf{nur auf ihren eigenen Adressraum}
            \item Kein direkter Zugriff auf Kernel-Speicher oder Hardware
            \item Müssen für Systemfunktionen einen \fatsf{System Call} machen $\rightarrow$ CPU wechselt in Kernel-Mode
        \end{itemize}
        \item \fatsf{Hardware Trap:} Eine \fatsf{Exception}, die direkt von der CPU ausgelöst wird, wenn etwas Unzulässiges passiert
        \item \fatsf{Segfault (Segmentation Fault):} Ein Laufzeitfehler, den das Betriebssystem an einen Prozess schickt, wenn dieser auf einen Speicherbereich zugreifen will, den er \fatsf{nicht darf}
    \end{itemize}
\end{infoBox}
\subsubsection*{Cache Side-Channel Angriffe}
\begin{figure}[h]
    \centering
    \includegraphics[width=1\textwidth]{Thema_8_Images/Cache Side-Channel Angriffe.jpeg}
\end{figure}
\begin{infoBox}
    \begin{itemize}
        \item \fatsf{Meltdown} ermöglicht es, \fatsf{geschützte Speicherbereiche} auszulesen, auf die ein normaler Prozess eigentlich \fatsf{keinen Zugriff} haben sollte
        \item \fatsf{Spectre} nutzt \fatsf{Spekulationsausführung}, um Daten aus anderen Prozessen zu stehlen
    \end{itemize}
\end{infoBox}
\clearpage
\subsubsection*{Cache Timing Side Channels}
\begin{itemize}
    \item Wir können nicht \fatsf{direkt} aus dem Cache lesen
    \item Grundidee:
    \begin{itemize}
        \item Wir können per Timing ermitteln, welche Adressen im Cache vorhanden sind
        \item Idee: Mache die Adressen \fatsf{abhängig} von den zu lesenden Daten
    \end{itemize}
    \item Viele verschiedene Techniken:
    \begin{itemize}
        \item Prime + Probe (Percival, 2005)
        \item Evict + Time (Osvik, Sheamir, Tromer, 2006)
        \item \fatsf{Flush + Reload (Yarom, Falkner, 2014)}
        \item Flush + Flush (Gruss, Maurice, Wagner, Mangard, 2016)
    \end{itemize}
\end{itemize}
\subsubsection*{Flush + Reload: Überblick}
\begin{infoBox}
    Flush + Reload ist ein \fatsf{Cache Side-Channel Angriff}, der es einem Angreifer ermöglicht, geheime Daten aus dem Speicher eines anderen Prozesses zu lesen, indem er die \fatsf{Cache-Zugriffszeiten} misst\\
    \fatsf{Ziel:}\\
    Den Wert eines \fatsf{geheimen Bytes} (data) herausfinden, indem überprüft wird, welche Speicheradresse schneller zugreifbar ist (also im Cache liegt)
\end{infoBox}
\fatsf{Ziel:} Lese ein Byte \blue{data} aus dem Speicher
\begin{enumerate}
    \item Allokiere Speicher für ein \red{probe\_array}, sorge dafür, dass keine Daten davon im Cache sind
    \item Lade einen Eintrag des Arrays in den Cache, dessen Adresse von \blue{data} abhängt
    \item Ermittle die im Cache vorhandene Adresse und berechne daraus \blue{data}
    \begin{itemize}
        \item \enquote{Kurzes} Timing = Adresse \fatsf{in} Cache vorhanden
        \item \enquote{Langes} Timing = Adresse \fatsf{nicht} im Cache
    \end{itemize}
\end{enumerate}
\clearpage
\subsubsection*{1. Vorbereitung des Probe Arrays}
\begin{figure}[h]
    \centering
    \includegraphics[width=0.8\textwidth]{Thema_8_Images/Vorbereitung des Probe Arrays.jpeg}
\end{figure}
\subsubsection*{Meltdown: Out-Of-Order Ausführung}
\begin{figure}[h]
    \centering
    \includegraphics[width=0.9\textwidth]{Thema_8_Images/Meltdown Out-of-order ausfürhung.jpeg}
\end{figure}
\clearpage
\subsubsection*{2. Array-Eintrag in den Cache laden}
\begin{figure}[h]
    \centering
    \includegraphics[width=1\textwidth]{Thema_8_Images/Array-Eintrag in dne Cache laden.jpeg}
\end{figure}
\subsubsection*{3. Data durch Timing ermitteln}
\begin{figure}[h]
    \centering
    \includegraphics[width=0.9\textwidth]{Thema_8_Images/Data durch Timing ermitteln 1.jpeg}
\end{figure}
\clearpage
\begin{figure}[h]
    \centering
    \includegraphics[width=1\textwidth]{Thema_8_Images/Data durch Timing ermitteln.jpeg}
\end{figure}
\subsubsection*{Meltdown}
\begin{itemize}
    \item Viele Intel CPUs seit $\approx$2010 verwundbar
    \begin{itemize}
        \item AMD, ARM: Schwieriger anzugreifen
    \end{itemize}
    \item Erlaubt das Lesen beliebiger Kernel-Adressen
    \begin{itemize}
        \item Physischer Speicher via Kernel-Mapping
    \end{itemize}
    \item Unabhängig von Software-Schwachstellen
    \begin{itemize}
        \item Basiert auf out-of-order Ausführung (\fatsf{Speculative Execution})
    \end{itemize}
\end{itemize}
\subsubsection*{Out-Of-Order Ausführung}
\begin{infoBox}
    Out-of-Order Execution (OoOE) ist eine Technik, die moderne Prozessorenverwenden, um die Leistung zu verbessern. Dabei werden Befehle \fatsf{nicht in der Reihenfolge des Programmcodes ausgeführt}, sondern so umsortiert, dass der Prozessor effizient arbeiten kann
\end{infoBox}
\begin{itemize}
    \item Ausführung zunächst ohne Zugriffskontrolle
    \begin{itemize}
        \item Kontrolle erst zum Ende der Ausführungspipeline
    \end{itemize}
    \item Instruktion \enquote{irrtümlich} ausgeführt:
    \begin{itemize}
        \item Rollback von Registern
        \item \fatsf{KEIN Cache-Rollback}
    \end{itemize}
\end{itemize}
\begin{infoBox}[title=Wie funktioniert Out-of-Order Execution?]
    \begin{enumerate}
        \item Der Prozessor analysiert die Befehle und sucht unabhängige Instruktionen
        \begin{itemize}
            \item Falls eine Instruktion auf das Ergebnis einer vorherigen wartet, kann eine spätere Instruktion bereits ausgeführt werden
        \end{itemize}
        \item Die Instruktionen werden \enquote{irrtümlich} ausgeführt, bevor alle Bedingungen geprüft werden
        \begin{itemize}
            \item Das bedeutet: Ein Zugriff auf geschützte Daten kann bereits stattfinden, bevor die CPU feststellt, dass der Zugriff nicht erlaubt ist
        \end{itemize}
        \item Am Ende der Pipeline wird geprüft, ob der Befehl erlaubt war
        \begin{itemize}
            \item Falls nicht, macht der Prozessor ein \fatsf{Rollback der Register} (die CPU \enquote{vergisst} die falsch Berechnung)
            \item \fatsf{Aber:} Der Cache bleibt unverändert - und genau hier liegt das Problem!
        \end{itemize}
    \end{enumerate}
\end{infoBox}
\subsubsection*{Meltdown mit Flush + Reload}
\begin{figure}[h]
    \centering
    \includegraphics[width=0.8\textwidth]{Thema_8_Images/Meltdown mit Flush-Reload.jpeg}
\end{figure}
\subsubsection*{Spectre}
\begin{itemize}
    \item High-Level Funktion wie Meltdown, aber Cache Side-Channel nun durch \fatsf{Branch Prediction}
\end{itemize}
\subsubsection*{Branch vorhersage}
\begin{figure}[h]
    \centering
    \includegraphics[width=1\textwidth]{Thema_8_Images/Branch vorhersage.jpeg}
\end{figure}
\begin{itemize}
    \item In den allermeisten Fällen wird \red{pos} innerhalb des Arrays liegen
    \item Moderne Prozessoren starten die Ausführung der nächsten Transaktion bevor die vorherige beendet ist \fatsf{(Pipelining)}
    \item \fatsf{Branch Prediction Unit} in der CPU entscheidet welche Verzweigung wahrscheinlicher ist. Diese wird \fatsf{spekulativ ausgeführt}
    \item Falsche Vorhersage führt zu Rollback, aber \fatsf{kein Cache Rollback}
    \item \fatsf{Branch Prediction Unit} \enquote{lernt} dynamisch um Vorhersage zu verbessern
\end{itemize}
\begin{infoBox}[title=Wie funktioniert Branch Prediction?]
    \begin{enumerate}
        \item Verzweigungen im Code (if/else, loops) verursachen Pipeline-Stalls:
        \begin{itemize}
            \item Beispiel:
            \begin{codeBlock}{}
                if (0 <= pos && pos < array_len){
                    res = array[pos];
                }
            \end{codeBlock}
            \item Die CPU \fatsf{weiß nicht sofort}, ob pos im gültigen Bereich ist, weil sie erst den Wert aus dem Speicher holen muss
        \end{itemize}
        \item Die CPU rät (spekulative Ausführung):
        \begin{itemize}
            \item Die \fatsf{Branch Prediction Unit} (BPU) sagt voraus, dass der häufigere Fall (pos < array\_len) eintreten wird
            \item Falls die Vorhersage falsch ist, wird ein \fatsf{Rollback durchgeführt}
            \item \fatsf{Problem:} Der Cache wird nicht zurückgesetzt, was zu Seitenkanal-Angriffen führen kann!
        \end{itemize}
    \end{enumerate}
\end{infoBox}
\begin{infoBox}[title=Missbrauch der Branch Prediction für Angriffe (Spectre)]
    Angreifer können die \fatsf{CPU trainieren}, um falsche Speicherzugriffe spekulativ auszuführen:\\
    \fatsf{Angriffsidee:}
    \begin{enumerate}
        \item Der Angreifer trainiert die CPU auf eine häufige Bedingung:
        \begin{itemize}
            \item Führt mehrmals \fatsf{valide} Zugriffe (pos < array\_len) aus, damit die CPU glaubt, dass das immer der Fall ist
        \end{itemize}
        \item Dann gibt er eine spezielle Eingabe, die außerhalb des erlaubten Bereichs liegt:
        \begin{itemize}
            \item pos = sehr große Zahl, sodass array1[pos] auf eine geschützte Kernel-Adresse zeigt
        \end{itemize}
        \item Die CPU glaubt, dass der Zugriff erlaubt ist und führt spekulativ aus:
        \begin{codeBlock}{}
            temp = probe_array[array1[x] * 4096];
        \end{codeBlock}
        \begin{itemize}
            \item Der Wert aus dem geschützten Speicher beeinflusst den Cache!
        \end{itemize}
        \item Der Angreifer misst den Cache und kann den geheimen Wert rekonstruieren!
    \end{enumerate}
\end{infoBox}
\begin{figure}[h]
    \centering
    \includegraphics[width=0.8\textwidth]{Thema_8_Images/Falsche Vorhersage der Branche.jpeg}
\end{figure}
\subsubsection*{Gegenmaßnahmen}
\begin{itemize}
    \item Einschränkung der spekulativen Ausführung
    \item Performanzeinbußen, Neu-Kompilierung von Anwendungen
    \item OpenBSD: Hyperthreading deaktiviert im Juni 2018 als Gegenmaßnahme für Spectre
\end{itemize}
$\Rightarrow$ Grundsätzliches Hardware/OS Re-Design

\subsubsection*{Lessons Learned}
\begin{itemize}
    \item Nur privilegierter Speicer im User-Adressraum $\rightarrow$ \blue{Sicher}
    \item Nur Out-of-Order und spekulative Ausfürhung $\rightarrow$ \blue{Sicher}
    \item Nur Caching $\rightarrow$ \blue{Sicher}
    \item Kombination $\rightarrow$ \red{Unsicher}
\end{itemize}
\begin{itemize}
    \item IT-Sicherheit muss das \enquote{Big Picture} im Blick haben
    \item Wir benötigen \fatsf{noch bessere Methodik}, um Security by Design zu erreichen!
\end{itemize}